\documentclass[11pt]{article}

%%%%%%%%%%%%%%Include Packages%%%%%%%%%%%%%%%%%%%%%%%%%%
\usepackage{xcolor}
\usepackage{mathtools}
\usepackage[a4paper, total={6in, 8in}, margin=1.25in]{geometry}
\usepackage{amsmath}
\usepackage{amssymb}
\usepackage{paralist}
\usepackage{rsfso}
\usepackage{amsthm}
\usepackage{wasysym}
\usepackage[inline]{enumitem}   
\usepackage{hyperref}
\usepackage{tocloft}
\usepackage{titlesec}
\usepackage{colortbl}
\usepackage{wrapfig}
\usepackage{pdfpages}
%%%%%%%%%%%%%%%%%%%%%%%%%%%%%%%%%%%%%%%%%%%%%%%%%%%%%%%%



\begin{document}
\Large \noindent \textbf{Physics 5A Midterm 1 Review Session}\\
\normalsize
Department of Physics and Astronomy, UCLA\\
\noindent\rule{8cm}{0.4pt}\\

\textbf{Problem 1:} During the discussion last week, Jinyan threw a stress ball against the wall. The ball left Jinyan's hand at $1$ meter above the ground with an initial velocity $10$ m/s at an angle $45^\circ$ above the horizon. The ball reach its maximum height before hitting the wall, and it hit the wall at a height $2$ meters above the ground.\\

\noindent\textbf{Part 1:} Find $x$- and $y$-components of the ball's initial velocity.\\
\textbf{Part 2:} Find the time it takes for the ball to hit the wall. \\
\textbf{Part 3:} Find the total horizontal distance it has traveled before reaching the wall.\\
\textbf{Part 4:} Find the maximal height the ball could reach.\\
\textbf{Part 5:} Find the ball's final velocity when it hit the wall (write it as $\vec{v} = (v_x,\,v_y)$).\\

\noindent Extra parts that you should know how to calculate:\\
\textbf{Part 6:} Find the total vertical displacement of the ball in its journey.\\
\textbf{Part 7:} Find the total vertical distance traveled in its journey.\\
\textbf{Part 8:} Find the total displacement vector and average velocity. 

\hfill\newline\hfill\newline\hfill\newline
\hfill\newline\hfill\newline\hfill\newline
\hfill\newline\hfill\newline\hfill\newline

\textbf{Problem 2}: Now instead of throwing the ball at an angle $45^\circ$ above the horizon, Jinyan throws the ball directly upward. The ball again starts at $1$ meter above the ground with an initial velocity $10$ m/s. The ball then drops to the ground after a few seconds. \\


\noindent\textbf{Part 1:} Find the time it takes for the ball to hit the ground. \\
\textbf{Part 2:} Find the maximal height the ball could reach. \\
\textbf{Part 3:} Find the total vertical distance it has traveled before hitting the ground.\\
\textbf{Part 4:} Find the average speed of the ball in its journey.\\
\textbf{Part 5:} Find the total displacement vector and average velocity.\\


\newpage
\includepdf[page=-]{5A_md1Rs}
 
\end{document}


